\section{Đề 16}
\debai{Trong mặt phẳng hệ tọa độ $Oxy$, cho đường thẳng $d: x-y+1-\sqrt{2}=0$ và điểm $A(-1;1)$. Viết phương trình đường tròn $(C)$ qua $A$, gốc tọa độ $O$ và tiếp xúc đường thẳng $d$. }

\debai{Giải hệ phương trình $\va{x^3+y^3+3(y-1)(x-y)=2 \\ \sqrt{x+1}+\sqrt{y+1}=\dfrac{(x-y)^2}{8}}$ }

\debai{Giả sử $x$ và $y$ không đồng thời bằng 0. Chứng minh
$$-2\sqrt{2}-2\le \dfrac{x^2-(x-4y)^2}{x^2+4y^2}\le 2\sqrt{2}-2$$}

\section{Đề 17}

\debai{Trong mặt phẳng tọa độ $Oxy$, cho tam giác $ABC$ nhọn có đỉnh $A(-1;4)$, trực tâm $H$. Đường thẳng $AH$ cắt cạnh $BC$ tại $M$, đường thẳng $CH$ cắt $AB$ tại $N$. Tâm đường tròn ngoại tiếp tam giác $HMN$ là $I(2;0)$, đường thẳng $BC$ đi qua điểm $P(1;-2)$. Tìm tọa độ các đỉnh $B,C$ của tam giác biết đỉnh $B$ thuộc đường thẳng $x+2y-2=0$.}

\debai{Giải hệ phương trình:
$ \va{
\dfrac{2}{(\sqrt{x}+\sqrt{y})^2}+\dfrac{1}{x+\sqrt{y(2x-y)}}=\dfrac{2}{y+\sqrt{x(2x-y)}}
\\

2(y-4)\sqrt{2x-y-3}-(x-6)\sqrt{x+y+1}=3(y-2)
 }} $
 
 \debai{Cho ba số thực $a,b,c$ thỏa mãn $a\geq2,b\geq 0,c\geq 0$. Tìm giá trị lớn nhất của biểu thức
 $$P=\dfrac{1}{2\sqrt{a^2+b^2+c^2-4a+5}}-\dfrac{1}{(a-1)(b+1)(c+1)}$$}
 
 \section{Đề 18}
 \debai{Trong mặt phẳng với hệ trục tọa độ $Oxy$. Viết phương trình các cạnh của hình vuông $ABCD$, biết rằng các đường thẳng $AB,CD,BC,AD$ lần lượt đi qua các điểm $M(2;4),N(2;-4),P(2;2),Q(3;-7)$.}
 
 \debai{Giải hệ phương trình:
$ \va{2x^2-y^2-7x+2y+6=0
 \\
 -7x^3+12x^2y-6xy^2+y^3-2x+2y=0 }
(x,y \in \mathbb{R}) $ }

\debai{Cho các số thực không âm $a,b,c$ thỏa mãn $a^2+b^2+c^2-3b \geq 0$. Tìm giá trị nhỏ nhất của biểu thức sau: 
$$ P=\dfrac{1}{(a+1)^2}+\dfrac{4}{(b+2)^2}+\dfrac{8}{(c+3)^2} $$}

\section{Đề 19}
\debai{Trong mặt phẳng với hệ trục tọa độ $Oxy$, cho hình vuông $ABCD$. Điểm $N(1;-2)$ thỏa mãn $2\vv{NB}+\vv{NC}=\vv{0}$ và điểm $M(3;6)$ thuộc đường thẳng chứa cạnh $AD$. Gọi $H$ là chân hình chiếu vuông góc của $A$ xuông đường thẳng $DN$. Xác định tọa độ của các đỉnh của hình vuông $ABCD$ biết khoảng cách từ điểm $H$ đến cạnh $CD$ bằng $\dfrac{12\sqrt{2}}{13}$ và đỉnh $A$ có hoành độ là một số nguyên lớn hơn -2}.

\debai{Giải hệ phương trình:
$\va{
\sqrt{x^2-x-y-1}.\sqrt[3]{x-y-1}=y+1 \\
x+y+1+\sqrt{2x+y}=\sqrt{5x^2+3y^2+3x+7y}
}$ $(x,y \in \mathbb{R})$}

\debai{Cho ba số thực không âm $x,y,z$. Tìm giá trị lớn nhất của biểu thức:
$$ P=\dfrac{4}{\sqrt{x^2+y^2+z^2+4}}-\dfrac{4}{(x+y)\sqrt{(x+2z)(y+2z)}}-\dfrac{5}{(y+z)\sqrt{(y+2x)(z+2x)}}$$}

\section{THPT Quỳnh Lưu 3 (Nghệ An)}
\debai{Trong mặt phẳng $Oxy$, cho hình chữ nhật $ABCD$ có $AB=2BC$. Gọi $H$ là hình chiếu của $A$ lên đường thẳng $BD;E,F$ lần lượt là trung điểm đoạn $CD, BH$. Biết $A(1;1)$, phương trình đường thẳng $BH$ là $3x-y-10=0$ và điểm $E$ có tung độ âm. Tìm tọa độ các đỉnh $B,C,D$.}

\debai{Giải hệ phương trình:
$\va{\sqrt{2(x+y+6)}=1-y \\
9\sqrt{1+x}+xy\sqrt{9+y^2}=0
}$}

\debai{Cho các số thực dương $a,b,c$ thỏa mãn $ab\geq1;c(a+b+c)\geq3$.
Tìm giá trị nhỏ nhất của biểu thức $$P=\dfrac{b+2c}{1+a}+\dfrac{a+2c}{1+b}+6\ln(a+b+2c)$$}

\section{THPT Thanh Chương III (Nghệ An)}
\debai{Trong mặt phẳng với hệ tọa độ $Oxy$ cho tam giác $ABC$ có $A(1;4)$, tiếp tuyến tại $A$ của đường tròn ngoại tiếp của tam giác $ABC$ cắt $BC$ tại $D$, đường phân giác trong của $\widehat{ADB}$ có phương trình $x-y+2=0$, điểm $M(-4;1)$ thuộc cạnh $AC$. Viết phương trình đường thẳng $AB$.}

\debai{Giải hệ phương trình:
$\va{x+\sqrt{xy+x-y^2-y}=5y+4 \\
\sqrt{4y^2-x-2}+\sqrt{y-1}=x-1
}$}

\debai{Cho $a,b,c$ là các số dương thỏa mãn $a+b+c=3$. Tìm giá trị lớn nhất của biểu thức: $$P=\dfrac{ab}{\sqrt{ab+3c}}+\dfrac{bc}{\sqrt{bc+3a}}+\dfrac{ca}{\sqrt{ca+3b}}$$}

\section{THPT Thiệu Hóa (Thanh Hóa)}
\debai{\begin{enumerate}
\item Trong mặt phẳng tọa độ $Oxy$, cho đường tròn $(C):x^2+y^2-4x+6y+4=0$. Viết phương trình các đường thẳng chứa các cạnh của hình vuông $MNPQ$ nội tiếp đường tròn $(C)$ biết điểm $M(2;0)$.
\item Trong mặt phẳng tọa độ $Oxy$, cho elip $(E):\dfrac{x^2}{16}+\dfrac{y^2}{9}=1.$ Tìm tọa độ các điểm $M$ trên $(E)$ sao cho $MF_{1}=2MF_{2} $ ( với $F_{1},F_{2}$ lần lượt là các tiêu điểm bên trái, bên phải của $(E)$).
\end{enumerate}
}
\debai{Giải hệ phương trình:
$\va{2.4^y+1=2^{1+\sqrt{2x}}+2\log_{2}{\dfrac{\sqrt{x}}{y}} \\
x^3+x=(y+1)(xy+1)+x^2
}$ $(x,y \in \mathbb{R})$.}

\debai{Cho $a,b,c$ là ba số thực dương. Chứng minh rằng:
$$\dfrac{a^2+1}{4b^2}+\dfrac{b^2+1}{4c^2}+\dfrac{c^2+1}{4a^2}\geq \dfrac{1}{a+b}+\dfrac{1}{b+c}+\dfrac{1}{c+a}.$$}

\section{THPT Thuận Châu (Sơn La)}
\debai{Trong mặt phẳng tọa độ $Oxy$, cho điểm $M(0;2)$ và hai đường thẳng $d:x+2y=0,\Delta :4x+3y=0$. Viết phương trình của đường tròn đi qua điểm $M$, có tâm thuộc đường thẳng $d$ và cắt đường thẳng $\Delta$ tại hai điểm phân biệt $A,B$ sao cho độ dài đoạn $AB$ bằng $4\sqrt{3}$. Biết tâm đường tròn có tung độ dương.}

\debai{Giải hệ phương trình:
$\va{x^3+12y^2+x+2=8y^3+8y \\
\sqrt{x^2+8y^3}=5x-2y
}$ $(x,y \in \mathbb{R}).$}

\debai{Tìm giá trị nhỏ nhất của biểu thức:
$$S=\dfrac{3}{-a+b+c}+\dfrac{4}{a-b+c}+\dfrac{5}{a+b-c}$$ Trong đó $a,b,c$ là độ dài của một tam giác thỏa mãn $2c+b=abc.$}

\section{THPT Tĩnh Gia I (Thanh Hóa)}
\debai{Trong mặt phẳng với hệ trục tọa độ $Oxy$, cho tam giác $ABC$ có trực tâm $H(3;0)$ và trung điểm của $BC$ là $I(6;1)$. Đường thẳng $AH$ có phương trình $x+2y-3=0$. Gọi $D,E$ lần lượt là chân đường cao kẻ từ $B$ và $C$ của tam giác $ABC$. Xác định tọa độ của cá đỉnh tam giác $ABC$, biết phương trình đường thẳng $DE$ là $x=2$ và điểm $D$ có tung độ dương.}

\debai{Giải hệ phương trình:
$\va{2y^2-3y+1+\sqrt{y-1}=x^2+\sqrt{x}+xy \\
\sqrt{2x+y}-\sqrt{-3x+2y+4}+3x^2-14x-8=0
}$}

\debai{Cho ba số thực không âm $a,b,c$ thỏa mãn $ab+bc+ca=1$. Chứng minh rằng: 
$$\dfrac{2a}{a^2+1}+\dfrac{2b}{b^2+1}+\dfrac{c^2-1}{c^2+1}\leqslant \dfrac{3}{2}.$$}

\section{THPT Thanh Chương I (Nghệ An)}
\debai{Trong mặt phẳng tọa độ $Oxy$, cho hình thang $ABCD$ có đường cao $AD$. Biết $BC=2AB, M(0;4)$ là trung điểm của $BC$ và phương trình đường thẳng $AD: x-2y-1=0$. Tìm tọa độ các đỉnh của hình thang biết diện tích hình thang là $\dfrac{54}{5}$ và $A,B$ có tọa độ dương.}

\debai{Giải hệ phương trình:
$\va{\sqrt{3y+1}+\sqrt{5x+4}=3xy-y+3 \\
\sqrt{2(x^2+y^2)}+\sqrt{\dfrac{4(x^2+xy+y^2)}{3}}=2(x+y)
}$ $(x,y \in \mathbb{R}).$}

\debai{Cho $a,b,c$ là các số thực dương thỏa mãn: $ab+bc+ca=3abc$. Tìm giá trị lớn nhất của biểu thức: $$A=\dfrac{1}{\sqrt{a^2+2b^2}}+\dfrac{1}{\sqrt{b^2+2c^2}}+\dfrac{1}{\sqrt{c^2+2a^2}}.$$} 

\section{THPT Cẩm Bình (Hà Tĩnh)}
\debai{Trong mặt phẳng với hệ tọa độ $Oxy$, cho tam giác $ABC$ có trực tâm $H(3;0)$. Biết $M(1;1),N(4;4)$ lần lượt là trung điểm của hai cạnh $AB,AC$. Tìm tọa độ các đỉnh của tam giác $ABC.$}

\debai{Giải hệ phương trình:
$\va{x^3-y^3+3y^2+32x=9x^2+8y+36 \\
4\sqrt{x+2}+\sqrt{16-3y}=x^2+8
}$ $(x,y \in \mathbb{R})$}

\debai{Cho $a,b,c$ là ba số thực dương. Tìm giá trị nhỏ nhất của biểu thức
$$P=\dfrac{a^2}{c(c^2+a^2)}+\dfrac{c^2}{b(b^2+c^2)}+\dfrac{b^2}{a(a^2+b^2)}.$$}

\section{THPT Lý Thái Tổ (Bắc Ninh)}
\debai{Trong mặt phẳng với hệ tọa độ $Oxy$, cho hình chữ nhật $ABCD$ có đường phân giác trong của góc $\widehat{ABC}$ đi qua trung điểm $M$ của cạnh $AD$, đường thẳng $BM$ có phương trình $x-y+2=0$, điểm $D$ nằm trên đường thẳng $\Delta: x+y-9=0$. Tìm tọa độ các đỉnh của hình chữ nhật $ABCD$ biết đỉnh $B$ có hoành độ âm và đường thẳng $AB$ đi qua $E(-1;2).$}

\debai{Giải hệ phương trình:
$\va{x^2-2x-2(x^2-x)\sqrt{3-2y}=(2y-3)x^2-1 \\
\sqrt{2-\sqrt{3-2y}}=\dfrac{\sqrt[3]{2x^2+x^3}+x+2}{2x+1}
}$}

\debai{Cho $x,y$ là hai số thỏa mãn: $x,y \geq1$ và $3(x+y)=4xy$. Tìm giá trị lớn nhất và giá trị nhỏ nhất của biểu thức:
$$P=x^3+y^3-3(\dfrac{1}{x^2}+\dfrac{1}{y^2}).$$}

\section{THPT Nghèn (Hà Tĩnh)}
\debai{Trong mặt phẳng tọa độ $Oxy$,cho hình vuông $ABCD$ có $M,N$ lần lượt là trung điểm $AB,BC$, biết $CM$ cắt $DN$ tại điểm $\left(\dfrac{22}{5};\dfrac{11}{5}\right)$. Gọi $H$ là trung điểm $DI$, biết đường thẳng $AH$ cắt $CD$ tại $P\left(\dfrac{7}{2};1\right)$. Tìm tọa độ các đỉnh của hình vuông $ABCD$ biết hoành độ điểm $A$ nhỏ hơn 4.}

\debai{Giải hệ phương trình:
$\va{(x^2+5y^2)^2=2\sqrt{xy}(6-x^2-5y^2)+36 \\
\sqrt{5y^4-x^4}=6x^2+2xy-6y^2
}$}

\debai{Cho $a,b,c$ là các số thực không đồng thời bằng 0 và thỏa mãn: $2(a^2+b^2+c^2)=(a+b+c)^2$. Tìm giá trị lớn nhất và giá trị nhỏ nhất của biểu thức:
$$P=\dfrac{a^3+b^3+c^3}{(a+b+c)(ab+bc+ca)}.$$}

\section{THPT chuyên Trần Quang Diệu (Đồng Tháp)}
\debai{Trong mặt phẳng với hệ tọa độ $Oxy$, cho đường tròn $(C):(x-2)^2+(y-2)^2=5$ và đường thẳng $\Delta: x+y+1=0$. Từ điểm $A$ thuộc đường thẳng $\Delta$, kẻ hai đường thẳng lần lượt tiếp xúc với $(C)$ tại $B$ và $C$. Tìm tọa độ điểm $A$ biết diện tích tam giác $ABC$ bằng 8.}

\debai{Giải hệ phương trình:
$\va{x^2y(2+2\sqrt{4y^2+1})=x+\sqrt{x^2+1} \\
x^2(4y^2+1)+2(x^2+1)\sqrt{x}=6
}$}

\debai{Cho các số thực không âm $a,b,c$ thỏa mãn $c=\min(a,b,c)$. Tìm giá trị nhỏ nhất của biểu thức:
$$P=\dfrac{1}{a^2+c^2}+\dfrac{1}{b^2+c^2}+\sqrt{a+b+c}.$$}

\section{THPT Nguyễn Thị Minh Khai (TPHCM)}
\debai{Trong mặt phẳng với hệ tọa độ $Oxy$ cho tam giác $ABC$ cân tại $B$, nội tiếp đường tròn $(C):x^2+y^2-10y-25=0$. $I$ là tâm đường tròn $(C)$. Đường thẳng $BI$ cắt đường tròn $(C)$ tại $M(5;0)$. Đường cao kẻ từ $C$ cắt đường tròn $(C)$ tại $N\left(\dfrac{-17}{5};\dfrac{-6}{5}\right)$. Tìm tạo độ các điểm $A,B,C$ biết điểm $A$ có hoành độ dương.}

\debai{Giải hệ phương trình:
$\va{x^3+3x^2+6x+4=y^3+3y \\
x^3(3y-7)=1-\sqrt{(1+x^2)^3}
}$ $(x,y \in \mathbb{R})$}

\debai{Cho các số dương $a,b,c$ thỏa mãn: $a(a-1)+b(b-1)+c(c-1)\leq \dfrac{4}{3}$}. Tìm giá trị nhỏ nhất của:
$$P=\dfrac{1}{a+1}+\dfrac{1}{b+1}+\dfrac{1}{c+1}$$